% =============================================================================
% CompetitiveExam/CS401/04-registers-instructions-answers.tex
% Answer key for 04-registers-instructions-questions.tex.
%
% Q1-Q7 are real GATE past-year questions, sourced from the local GATEOverflow
% Volume 3 PYQ book (CompetitiveExam/Books/index.md, Tier-1 availability).
% The book authenticates each question (year + question number printed on the
% page; full text recovered via render+OCR). The answers below are derived and
% justified step by step. Each derivation is self-contained so the answers can
% be independently re-checked; the collapsed answer keys inside the book are
% not extractable, so these justifications stand in for the book's key.
% =============================================================================

\documentclass[11pt]{article}
\usepackage[margin=1in]{geometry}
% =============================================================================
% Preambles/header.tex — shared LaTeX/Beamer preamble for this project
%
% Use from a Beamer lecture by prepending:
%     % =============================================================================
% Preambles/header.tex — shared LaTeX/Beamer preamble for this project
%
% Use from a Beamer lecture by prepending:
%     % =============================================================================
% Preambles/header.tex — shared LaTeX/Beamer preamble for this project
%
% Use from a Beamer lecture by prepending:
%     \input{header}
% and compiling with TEXINPUTS=../Preambles:$TEXINPUTS (the /compile-latex
% skill does this automatically).
%
% PALETTE CONTRACT. Color names defined here MUST match the names in
% Quarto/theme-template.scss so Beamer and Quarto renderings use the same
% palette. The scripts/check-palette-sync.sh script enforces this.
%
% To customize: edit the HEX values below AND the matching $variable in
% Quarto/theme-template.scss, then run scripts/check-palette-sync.sh.
%
% TITLE PAGE. A formal, serif (Libertine) title-page template is defined in
% the Beamer-only block below (\setbeamertemplate{title page}) -- title,
% subtitle, a thin gold rule, then \author/\institute/\date. Frametitle and
% body text remain on Beamer's sans-serif default. No SCSS counterpart is
% needed for this (font/layout only, no new colors).
% =============================================================================

% -----------------------------------------------------------------------------
% Engine detection + encoding
%
% Our /compile-latex skill uses XeLaTeX, where inputenc is unnecessary and
% can produce encoding warnings. Only load inputenc under pdfTeX.
% -----------------------------------------------------------------------------
\usepackage{iftex}
\ifPDFTeX
  \usepackage[utf8]{inputenc}
\fi

\usepackage{amsmath, amssymb, amsthm}
\usepackage{booktabs}
\usepackage{graphicx}

% -----------------------------------------------------------------------------
% Serif face for the title page (formal/academic look). Linux Libertine is
% XeLaTeX- and pdfTeX-safe and redefines \rmdefault, so \rmfamily picks it up
% everywhere \rmfamily is used explicitly. Body text and frametitles stay on
% Beamer's own sans-serif default (unaffected) -- only the title-page
% \setbeamerfont{...} overrides below opt in to \rmfamily.
% -----------------------------------------------------------------------------
\usepackage{libertine}

% -----------------------------------------------------------------------------
% PALETTE — mirrors Quarto/theme-template.scss variable names.
% Load xcolor and define colors BEFORE anything that references them
% (notably hyperref, hypersetup, and the Beamer color assignments).
% -----------------------------------------------------------------------------
\usepackage{xcolor}

\definecolor{primary-blue}{HTML}{012169}      % headings, accents
\definecolor{primary-gold}{HTML}{B9975B}      % emphasis, borders
\definecolor{highlight-yellow}{HTML}{F2A900}  % markers, alerts
\definecolor{light-bg}{HTML}{E8EDF5}          % title / section backgrounds
\definecolor{jet}{HTML}{1A1A1A}               % body text

% Semantic colors (used in diagrams and annotations)
\definecolor{positive}{HTML}{15803D}          % good / observed / identified
\definecolor{negative}{HTML}{B91C1C}          % bad / problematic / bias
\definecolor{neutral}{HTML}{525252}           % reference / context

% Highlight accents (match .hi-* classes in the SCSS)
\definecolor{hi-slate}{HTML}{314F4F}
\definecolor{hi-green}{HTML}{2E8B57}
\definecolor{hi-red}{HTML}{C41E3A}

% Semantic aliases so skill templates and snippets can write
% \textcolor{accent}{...} without hard-coding "primary-blue".
\colorlet{accent}{primary-blue}
\colorlet{accent2}{primary-gold}

% -----------------------------------------------------------------------------
% Hyperref — loaded AFTER xcolor + palette so link colors resolve.
% -----------------------------------------------------------------------------
\usepackage{hyperref}
\hypersetup{colorlinks=true, linkcolor=primary-blue, urlcolor=primary-blue,
            citecolor=primary-blue, pdfborder={0 0 0}}

% -----------------------------------------------------------------------------
% Course code — set once per deck (after \input{header}, before \begin{document})
% via \coursecode{CS401}. Shown in the footer on every slide so a deck's course
% is identifiable without opening the title page. Empty by default.
% -----------------------------------------------------------------------------
\makeatletter
\newcommand{\coursecode}[1]{\gdef\@coursecodetext{#1}}
\gdef\@coursecodetext{}
\makeatother

% -----------------------------------------------------------------------------
% Beamer theme assignments (only apply under Beamer).
%
% We detect Beamer via \@ifclassloaded{beamer}, which is the documented,
% non-internal way. This must be wrapped in \makeatletter/\makeatother
% because \@ifclassloaded contains an @-catcode character.
% -----------------------------------------------------------------------------
\makeatletter
\@ifclassloaded{beamer}{%
  \usefonttheme{professionalfonts}
  \setbeamercolor{structure}{fg=primary-blue}
  \setbeamercolor{frametitle}{fg=primary-blue}
  \setbeamercolor{title}{fg=primary-blue}
  \setbeamercolor{subtitle}{fg=primary-gold}
  \setbeamercolor{author}{fg=primary-blue}
  \setbeamercolor{institute}{fg=neutral}
  \setbeamercolor{date}{fg=primary-gold}
  \setbeamercolor{itemize item}{fg=primary-blue}
  \setbeamercolor{itemize subitem}{fg=primary-gold}
  \setbeamercolor{alerted text}{fg=primary-gold}
  \setbeamercolor{block title}{fg=white, bg=primary-blue}
  \setbeamercolor{block body}{bg=light-bg}
  % Semantic block variants: block=definition/concept (blue), exampleblock=
  % worked example (green/positive), alertblock=key takeaway or caution (gold).
  % Keep to <=2 colored boxes per slide across ALL three variants (INV-7).
  \setbeamercolor{block title example}{fg=white, bg=positive}
  \setbeamercolor{block body example}{bg=light-bg}
  \setbeamercolor{block title alerted}{fg=white, bg=primary-gold}
  \setbeamercolor{block body alerted}{bg=light-bg}

  % ---------------------------------------------------------------------------
  % Formal title page: serif (Libertine, via \rmfamily) for title-page
  % elements only. Body text, itemize, and frametitle are intentionally left
  % on Beamer's sans-serif default -- \usebeamerfont{title} is also what
  % \transitionslide (below) uses for its big centered text, so section
  % transitions pick up the same serif treatment for free.
  % ---------------------------------------------------------------------------
  \setbeamerfont{title}{family=\rmfamily, series=\bfseries, size=\Large}
  \setbeamerfont{subtitle}{family=\rmfamily, shape=\itshape, size=\normalsize}
  \setbeamerfont{author}{family=\rmfamily, size=\normalsize}
  \setbeamerfont{institute}{family=\rmfamily, size=\scriptsize}
  \setbeamerfont{date}{family=\rmfamily, size=\scriptsize}

  \setbeamertemplate{title page}{
    \vfill
    \begin{center}
      {\usebeamerfont{title}\usebeamercolor[fg]{title}\inserttitle\par}
      \vspace{0.15cm}
      {\usebeamerfont{subtitle}\usebeamercolor[fg]{subtitle}\insertsubtitle\par}
      \vspace{0.45cm}
      {\color{primary-gold}\rule{0.42\textwidth}{0.6pt}\par}
      \vspace{0.45cm}
      {\usebeamerfont{author}\usebeamercolor[fg]{author}\insertauthor\par}
      \vspace{0.35cm}
      {\usebeamerfont{institute}\usebeamercolor[fg]{institute}\insertinstitute\par}
      \vspace{0.35cm}
      {\usebeamerfont{date}\usebeamercolor[fg]{date}\insertdate\par}
    \end{center}
    \vfill
  }

  % Minimal footer: course code (left) + page X / N (right), no navigation clutter
  \setbeamertemplate{navigation symbols}{}
  \setbeamertemplate{footline}{%
    \color{neutral}\footnotesize\hspace{0.7em}\@coursecodetext\hfill
    \insertframenumber/\inserttotalframenumber\hspace{0.7em}\vspace{0.3em}}
}{}
\makeatother

% -----------------------------------------------------------------------------
% TikZ setup — libraries the snippet gallery uses
% -----------------------------------------------------------------------------
\usepackage{tikz}
\usetikzlibrary{arrows.meta, positioning, calc, decorations.pathreplacing,
                fit, shapes.geometric, backgrounds}

% Shared TikZ styles — snippets and hand-written diagrams can pick these up
% via `\begin{tikzpicture}[every node/.style={...}]` or by naming specific
% styles (e.g., `\node[dag-node] (X) at (0,0) {X};`).
\tikzset{
  % Node styles for causal DAGs and similar
  dag-node/.style = {
    draw=primary-blue, fill=light-bg, rounded corners=2pt,
    minimum width=1.2cm, minimum height=0.9cm, align=center, font=\small
  },
  decision-node/.style = {
    draw=primary-gold, fill=white, diamond, aspect=1.8,
    minimum width=1.8cm, minimum height=1.2cm, align=center, font=\small,
    inner sep=1pt
  },
  % Edge styles — observed vs counterfactual
  observed-edge/.style = {-{Latex[length=2.5mm]}, thick, draw=jet},
  counterfactual-edge/.style = {-{Latex[length=2.5mm]}, thick, draw=neutral, dashed},
  confound-edge/.style = {-{Latex[length=2.5mm]}, thick, draw=negative, dashed},
  % Data-point styles for plots
  observed-dot/.style = {fill=primary-blue, draw=primary-blue, circle, inner sep=1.2pt},
  counterfactual-dot/.style = {fill=white, draw=neutral, circle, inner sep=1.2pt}
}

% -----------------------------------------------------------------------------
% Convenience macros
% -----------------------------------------------------------------------------
\newcommand{\muted}[1]{\textcolor{neutral}{#1}}
\newcommand{\key}[1]{\textcolor{primary-gold}{\textbf{#1}}}
\newcommand{\good}[1]{\textcolor{positive}{#1}}
\newcommand{\bad}[1]{\textcolor{negative}{#1}}

% Standout frame macro: full-bleed dark background for section transitions.
% Beamer-only (uses \begin{frame}/\setbeamercolor/\usebeamerfont) -- do not
% call from an article-class file (e.g. Notes/<CODE>/*-notes.tex). \muted,
% \key, \good, \bad, and \coursecode above are plain xcolor/gdef and are
% safe to use from either class; verified when Notes/article-class support
% was added.
% Expand: \transitionslide{Your transition title}
\newcommand{\transitionslide}[1]{%
  \begingroup
  \setbeamercolor{background canvas}{bg=primary-blue}
  \begin{frame}[plain]
    \centering
    \vfill
    {\usebeamerfont{title}\color{white}\Large #1\par}
    \vfill
  \end{frame}
  \endgroup
}

% and compiling with TEXINPUTS=../Preambles:$TEXINPUTS (the /compile-latex
% skill does this automatically).
%
% PALETTE CONTRACT. Color names defined here MUST match the names in
% Quarto/theme-template.scss so Beamer and Quarto renderings use the same
% palette. The scripts/check-palette-sync.sh script enforces this.
%
% To customize: edit the HEX values below AND the matching $variable in
% Quarto/theme-template.scss, then run scripts/check-palette-sync.sh.
%
% TITLE PAGE. A formal, serif (Libertine) title-page template is defined in
% the Beamer-only block below (\setbeamertemplate{title page}) -- title,
% subtitle, a thin gold rule, then \author/\institute/\date. Frametitle and
% body text remain on Beamer's sans-serif default. No SCSS counterpart is
% needed for this (font/layout only, no new colors).
% =============================================================================

% -----------------------------------------------------------------------------
% Engine detection + encoding
%
% Our /compile-latex skill uses XeLaTeX, where inputenc is unnecessary and
% can produce encoding warnings. Only load inputenc under pdfTeX.
% -----------------------------------------------------------------------------
\usepackage{iftex}
\ifPDFTeX
  \usepackage[utf8]{inputenc}
\fi

\usepackage{amsmath, amssymb, amsthm}
\usepackage{booktabs}
\usepackage{graphicx}

% -----------------------------------------------------------------------------
% Serif face for the title page (formal/academic look). Linux Libertine is
% XeLaTeX- and pdfTeX-safe and redefines \rmdefault, so \rmfamily picks it up
% everywhere \rmfamily is used explicitly. Body text and frametitles stay on
% Beamer's own sans-serif default (unaffected) -- only the title-page
% \setbeamerfont{...} overrides below opt in to \rmfamily.
% -----------------------------------------------------------------------------
\usepackage{libertine}

% -----------------------------------------------------------------------------
% PALETTE — mirrors Quarto/theme-template.scss variable names.
% Load xcolor and define colors BEFORE anything that references them
% (notably hyperref, hypersetup, and the Beamer color assignments).
% -----------------------------------------------------------------------------
\usepackage{xcolor}

\definecolor{primary-blue}{HTML}{012169}      % headings, accents
\definecolor{primary-gold}{HTML}{B9975B}      % emphasis, borders
\definecolor{highlight-yellow}{HTML}{F2A900}  % markers, alerts
\definecolor{light-bg}{HTML}{E8EDF5}          % title / section backgrounds
\definecolor{jet}{HTML}{1A1A1A}               % body text

% Semantic colors (used in diagrams and annotations)
\definecolor{positive}{HTML}{15803D}          % good / observed / identified
\definecolor{negative}{HTML}{B91C1C}          % bad / problematic / bias
\definecolor{neutral}{HTML}{525252}           % reference / context

% Highlight accents (match .hi-* classes in the SCSS)
\definecolor{hi-slate}{HTML}{314F4F}
\definecolor{hi-green}{HTML}{2E8B57}
\definecolor{hi-red}{HTML}{C41E3A}

% Semantic aliases so skill templates and snippets can write
% \textcolor{accent}{...} without hard-coding "primary-blue".
\colorlet{accent}{primary-blue}
\colorlet{accent2}{primary-gold}

% -----------------------------------------------------------------------------
% Hyperref — loaded AFTER xcolor + palette so link colors resolve.
% -----------------------------------------------------------------------------
\usepackage{hyperref}
\hypersetup{colorlinks=true, linkcolor=primary-blue, urlcolor=primary-blue,
            citecolor=primary-blue, pdfborder={0 0 0}}

% -----------------------------------------------------------------------------
% Course code — set once per deck (after % =============================================================================
% Preambles/header.tex — shared LaTeX/Beamer preamble for this project
%
% Use from a Beamer lecture by prepending:
%     \input{header}
% and compiling with TEXINPUTS=../Preambles:$TEXINPUTS (the /compile-latex
% skill does this automatically).
%
% PALETTE CONTRACT. Color names defined here MUST match the names in
% Quarto/theme-template.scss so Beamer and Quarto renderings use the same
% palette. The scripts/check-palette-sync.sh script enforces this.
%
% To customize: edit the HEX values below AND the matching $variable in
% Quarto/theme-template.scss, then run scripts/check-palette-sync.sh.
%
% TITLE PAGE. A formal, serif (Libertine) title-page template is defined in
% the Beamer-only block below (\setbeamertemplate{title page}) -- title,
% subtitle, a thin gold rule, then \author/\institute/\date. Frametitle and
% body text remain on Beamer's sans-serif default. No SCSS counterpart is
% needed for this (font/layout only, no new colors).
% =============================================================================

% -----------------------------------------------------------------------------
% Engine detection + encoding
%
% Our /compile-latex skill uses XeLaTeX, where inputenc is unnecessary and
% can produce encoding warnings. Only load inputenc under pdfTeX.
% -----------------------------------------------------------------------------
\usepackage{iftex}
\ifPDFTeX
  \usepackage[utf8]{inputenc}
\fi

\usepackage{amsmath, amssymb, amsthm}
\usepackage{booktabs}
\usepackage{graphicx}

% -----------------------------------------------------------------------------
% Serif face for the title page (formal/academic look). Linux Libertine is
% XeLaTeX- and pdfTeX-safe and redefines \rmdefault, so \rmfamily picks it up
% everywhere \rmfamily is used explicitly. Body text and frametitles stay on
% Beamer's own sans-serif default (unaffected) -- only the title-page
% \setbeamerfont{...} overrides below opt in to \rmfamily.
% -----------------------------------------------------------------------------
\usepackage{libertine}

% -----------------------------------------------------------------------------
% PALETTE — mirrors Quarto/theme-template.scss variable names.
% Load xcolor and define colors BEFORE anything that references them
% (notably hyperref, hypersetup, and the Beamer color assignments).
% -----------------------------------------------------------------------------
\usepackage{xcolor}

\definecolor{primary-blue}{HTML}{012169}      % headings, accents
\definecolor{primary-gold}{HTML}{B9975B}      % emphasis, borders
\definecolor{highlight-yellow}{HTML}{F2A900}  % markers, alerts
\definecolor{light-bg}{HTML}{E8EDF5}          % title / section backgrounds
\definecolor{jet}{HTML}{1A1A1A}               % body text

% Semantic colors (used in diagrams and annotations)
\definecolor{positive}{HTML}{15803D}          % good / observed / identified
\definecolor{negative}{HTML}{B91C1C}          % bad / problematic / bias
\definecolor{neutral}{HTML}{525252}           % reference / context

% Highlight accents (match .hi-* classes in the SCSS)
\definecolor{hi-slate}{HTML}{314F4F}
\definecolor{hi-green}{HTML}{2E8B57}
\definecolor{hi-red}{HTML}{C41E3A}

% Semantic aliases so skill templates and snippets can write
% \textcolor{accent}{...} without hard-coding "primary-blue".
\colorlet{accent}{primary-blue}
\colorlet{accent2}{primary-gold}

% -----------------------------------------------------------------------------
% Hyperref — loaded AFTER xcolor + palette so link colors resolve.
% -----------------------------------------------------------------------------
\usepackage{hyperref}
\hypersetup{colorlinks=true, linkcolor=primary-blue, urlcolor=primary-blue,
            citecolor=primary-blue, pdfborder={0 0 0}}

% -----------------------------------------------------------------------------
% Course code — set once per deck (after \input{header}, before \begin{document})
% via \coursecode{CS401}. Shown in the footer on every slide so a deck's course
% is identifiable without opening the title page. Empty by default.
% -----------------------------------------------------------------------------
\makeatletter
\newcommand{\coursecode}[1]{\gdef\@coursecodetext{#1}}
\gdef\@coursecodetext{}
\makeatother

% -----------------------------------------------------------------------------
% Beamer theme assignments (only apply under Beamer).
%
% We detect Beamer via \@ifclassloaded{beamer}, which is the documented,
% non-internal way. This must be wrapped in \makeatletter/\makeatother
% because \@ifclassloaded contains an @-catcode character.
% -----------------------------------------------------------------------------
\makeatletter
\@ifclassloaded{beamer}{%
  \usefonttheme{professionalfonts}
  \setbeamercolor{structure}{fg=primary-blue}
  \setbeamercolor{frametitle}{fg=primary-blue}
  \setbeamercolor{title}{fg=primary-blue}
  \setbeamercolor{subtitle}{fg=primary-gold}
  \setbeamercolor{author}{fg=primary-blue}
  \setbeamercolor{institute}{fg=neutral}
  \setbeamercolor{date}{fg=primary-gold}
  \setbeamercolor{itemize item}{fg=primary-blue}
  \setbeamercolor{itemize subitem}{fg=primary-gold}
  \setbeamercolor{alerted text}{fg=primary-gold}
  \setbeamercolor{block title}{fg=white, bg=primary-blue}
  \setbeamercolor{block body}{bg=light-bg}
  % Semantic block variants: block=definition/concept (blue), exampleblock=
  % worked example (green/positive), alertblock=key takeaway or caution (gold).
  % Keep to <=2 colored boxes per slide across ALL three variants (INV-7).
  \setbeamercolor{block title example}{fg=white, bg=positive}
  \setbeamercolor{block body example}{bg=light-bg}
  \setbeamercolor{block title alerted}{fg=white, bg=primary-gold}
  \setbeamercolor{block body alerted}{bg=light-bg}

  % ---------------------------------------------------------------------------
  % Formal title page: serif (Libertine, via \rmfamily) for title-page
  % elements only. Body text, itemize, and frametitle are intentionally left
  % on Beamer's sans-serif default -- \usebeamerfont{title} is also what
  % \transitionslide (below) uses for its big centered text, so section
  % transitions pick up the same serif treatment for free.
  % ---------------------------------------------------------------------------
  \setbeamerfont{title}{family=\rmfamily, series=\bfseries, size=\Large}
  \setbeamerfont{subtitle}{family=\rmfamily, shape=\itshape, size=\normalsize}
  \setbeamerfont{author}{family=\rmfamily, size=\normalsize}
  \setbeamerfont{institute}{family=\rmfamily, size=\scriptsize}
  \setbeamerfont{date}{family=\rmfamily, size=\scriptsize}

  \setbeamertemplate{title page}{
    \vfill
    \begin{center}
      {\usebeamerfont{title}\usebeamercolor[fg]{title}\inserttitle\par}
      \vspace{0.15cm}
      {\usebeamerfont{subtitle}\usebeamercolor[fg]{subtitle}\insertsubtitle\par}
      \vspace{0.45cm}
      {\color{primary-gold}\rule{0.42\textwidth}{0.6pt}\par}
      \vspace{0.45cm}
      {\usebeamerfont{author}\usebeamercolor[fg]{author}\insertauthor\par}
      \vspace{0.35cm}
      {\usebeamerfont{institute}\usebeamercolor[fg]{institute}\insertinstitute\par}
      \vspace{0.35cm}
      {\usebeamerfont{date}\usebeamercolor[fg]{date}\insertdate\par}
    \end{center}
    \vfill
  }

  % Minimal footer: course code (left) + page X / N (right), no navigation clutter
  \setbeamertemplate{navigation symbols}{}
  \setbeamertemplate{footline}{%
    \color{neutral}\footnotesize\hspace{0.7em}\@coursecodetext\hfill
    \insertframenumber/\inserttotalframenumber\hspace{0.7em}\vspace{0.3em}}
}{}
\makeatother

% -----------------------------------------------------------------------------
% TikZ setup — libraries the snippet gallery uses
% -----------------------------------------------------------------------------
\usepackage{tikz}
\usetikzlibrary{arrows.meta, positioning, calc, decorations.pathreplacing,
                fit, shapes.geometric, backgrounds}

% Shared TikZ styles — snippets and hand-written diagrams can pick these up
% via `\begin{tikzpicture}[every node/.style={...}]` or by naming specific
% styles (e.g., `\node[dag-node] (X) at (0,0) {X};`).
\tikzset{
  % Node styles for causal DAGs and similar
  dag-node/.style = {
    draw=primary-blue, fill=light-bg, rounded corners=2pt,
    minimum width=1.2cm, minimum height=0.9cm, align=center, font=\small
  },
  decision-node/.style = {
    draw=primary-gold, fill=white, diamond, aspect=1.8,
    minimum width=1.8cm, minimum height=1.2cm, align=center, font=\small,
    inner sep=1pt
  },
  % Edge styles — observed vs counterfactual
  observed-edge/.style = {-{Latex[length=2.5mm]}, thick, draw=jet},
  counterfactual-edge/.style = {-{Latex[length=2.5mm]}, thick, draw=neutral, dashed},
  confound-edge/.style = {-{Latex[length=2.5mm]}, thick, draw=negative, dashed},
  % Data-point styles for plots
  observed-dot/.style = {fill=primary-blue, draw=primary-blue, circle, inner sep=1.2pt},
  counterfactual-dot/.style = {fill=white, draw=neutral, circle, inner sep=1.2pt}
}

% -----------------------------------------------------------------------------
% Convenience macros
% -----------------------------------------------------------------------------
\newcommand{\muted}[1]{\textcolor{neutral}{#1}}
\newcommand{\key}[1]{\textcolor{primary-gold}{\textbf{#1}}}
\newcommand{\good}[1]{\textcolor{positive}{#1}}
\newcommand{\bad}[1]{\textcolor{negative}{#1}}

% Standout frame macro: full-bleed dark background for section transitions.
% Beamer-only (uses \begin{frame}/\setbeamercolor/\usebeamerfont) -- do not
% call from an article-class file (e.g. Notes/<CODE>/*-notes.tex). \muted,
% \key, \good, \bad, and \coursecode above are plain xcolor/gdef and are
% safe to use from either class; verified when Notes/article-class support
% was added.
% Expand: \transitionslide{Your transition title}
\newcommand{\transitionslide}[1]{%
  \begingroup
  \setbeamercolor{background canvas}{bg=primary-blue}
  \begin{frame}[plain]
    \centering
    \vfill
    {\usebeamerfont{title}\color{white}\Large #1\par}
    \vfill
  \end{frame}
  \endgroup
}
, before \begin{document})
% via \coursecode{CS401}. Shown in the footer on every slide so a deck's course
% is identifiable without opening the title page. Empty by default.
% -----------------------------------------------------------------------------
\makeatletter
\newcommand{\coursecode}[1]{\gdef\@coursecodetext{#1}}
\gdef\@coursecodetext{}
\makeatother

% -----------------------------------------------------------------------------
% Beamer theme assignments (only apply under Beamer).
%
% We detect Beamer via \@ifclassloaded{beamer}, which is the documented,
% non-internal way. This must be wrapped in \makeatletter/\makeatother
% because \@ifclassloaded contains an @-catcode character.
% -----------------------------------------------------------------------------
\makeatletter
\@ifclassloaded{beamer}{%
  \usefonttheme{professionalfonts}
  \setbeamercolor{structure}{fg=primary-blue}
  \setbeamercolor{frametitle}{fg=primary-blue}
  \setbeamercolor{title}{fg=primary-blue}
  \setbeamercolor{subtitle}{fg=primary-gold}
  \setbeamercolor{author}{fg=primary-blue}
  \setbeamercolor{institute}{fg=neutral}
  \setbeamercolor{date}{fg=primary-gold}
  \setbeamercolor{itemize item}{fg=primary-blue}
  \setbeamercolor{itemize subitem}{fg=primary-gold}
  \setbeamercolor{alerted text}{fg=primary-gold}
  \setbeamercolor{block title}{fg=white, bg=primary-blue}
  \setbeamercolor{block body}{bg=light-bg}
  % Semantic block variants: block=definition/concept (blue), exampleblock=
  % worked example (green/positive), alertblock=key takeaway or caution (gold).
  % Keep to <=2 colored boxes per slide across ALL three variants (INV-7).
  \setbeamercolor{block title example}{fg=white, bg=positive}
  \setbeamercolor{block body example}{bg=light-bg}
  \setbeamercolor{block title alerted}{fg=white, bg=primary-gold}
  \setbeamercolor{block body alerted}{bg=light-bg}

  % ---------------------------------------------------------------------------
  % Formal title page: serif (Libertine, via \rmfamily) for title-page
  % elements only. Body text, itemize, and frametitle are intentionally left
  % on Beamer's sans-serif default -- \usebeamerfont{title} is also what
  % \transitionslide (below) uses for its big centered text, so section
  % transitions pick up the same serif treatment for free.
  % ---------------------------------------------------------------------------
  \setbeamerfont{title}{family=\rmfamily, series=\bfseries, size=\Large}
  \setbeamerfont{subtitle}{family=\rmfamily, shape=\itshape, size=\normalsize}
  \setbeamerfont{author}{family=\rmfamily, size=\normalsize}
  \setbeamerfont{institute}{family=\rmfamily, size=\scriptsize}
  \setbeamerfont{date}{family=\rmfamily, size=\scriptsize}

  \setbeamertemplate{title page}{
    \vfill
    \begin{center}
      {\usebeamerfont{title}\usebeamercolor[fg]{title}\inserttitle\par}
      \vspace{0.15cm}
      {\usebeamerfont{subtitle}\usebeamercolor[fg]{subtitle}\insertsubtitle\par}
      \vspace{0.45cm}
      {\color{primary-gold}\rule{0.42\textwidth}{0.6pt}\par}
      \vspace{0.45cm}
      {\usebeamerfont{author}\usebeamercolor[fg]{author}\insertauthor\par}
      \vspace{0.35cm}
      {\usebeamerfont{institute}\usebeamercolor[fg]{institute}\insertinstitute\par}
      \vspace{0.35cm}
      {\usebeamerfont{date}\usebeamercolor[fg]{date}\insertdate\par}
    \end{center}
    \vfill
  }

  % Minimal footer: course code (left) + page X / N (right), no navigation clutter
  \setbeamertemplate{navigation symbols}{}
  \setbeamertemplate{footline}{%
    \color{neutral}\footnotesize\hspace{0.7em}\@coursecodetext\hfill
    \insertframenumber/\inserttotalframenumber\hspace{0.7em}\vspace{0.3em}}
}{}
\makeatother

% -----------------------------------------------------------------------------
% TikZ setup — libraries the snippet gallery uses
% -----------------------------------------------------------------------------
\usepackage{tikz}
\usetikzlibrary{arrows.meta, positioning, calc, decorations.pathreplacing,
                fit, shapes.geometric, backgrounds}

% Shared TikZ styles — snippets and hand-written diagrams can pick these up
% via `\begin{tikzpicture}[every node/.style={...}]` or by naming specific
% styles (e.g., `\node[dag-node] (X) at (0,0) {X};`).
\tikzset{
  % Node styles for causal DAGs and similar
  dag-node/.style = {
    draw=primary-blue, fill=light-bg, rounded corners=2pt,
    minimum width=1.2cm, minimum height=0.9cm, align=center, font=\small
  },
  decision-node/.style = {
    draw=primary-gold, fill=white, diamond, aspect=1.8,
    minimum width=1.8cm, minimum height=1.2cm, align=center, font=\small,
    inner sep=1pt
  },
  % Edge styles — observed vs counterfactual
  observed-edge/.style = {-{Latex[length=2.5mm]}, thick, draw=jet},
  counterfactual-edge/.style = {-{Latex[length=2.5mm]}, thick, draw=neutral, dashed},
  confound-edge/.style = {-{Latex[length=2.5mm]}, thick, draw=negative, dashed},
  % Data-point styles for plots
  observed-dot/.style = {fill=primary-blue, draw=primary-blue, circle, inner sep=1.2pt},
  counterfactual-dot/.style = {fill=white, draw=neutral, circle, inner sep=1.2pt}
}

% -----------------------------------------------------------------------------
% Convenience macros
% -----------------------------------------------------------------------------
\newcommand{\muted}[1]{\textcolor{neutral}{#1}}
\newcommand{\key}[1]{\textcolor{primary-gold}{\textbf{#1}}}
\newcommand{\good}[1]{\textcolor{positive}{#1}}
\newcommand{\bad}[1]{\textcolor{negative}{#1}}

% Standout frame macro: full-bleed dark background for section transitions.
% Beamer-only (uses \begin{frame}/\setbeamercolor/\usebeamerfont) -- do not
% call from an article-class file (e.g. Notes/<CODE>/*-notes.tex). \muted,
% \key, \good, \bad, and \coursecode above are plain xcolor/gdef and are
% safe to use from either class; verified when Notes/article-class support
% was added.
% Expand: \transitionslide{Your transition title}
\newcommand{\transitionslide}[1]{%
  \begingroup
  \setbeamercolor{background canvas}{bg=primary-blue}
  \begin{frame}[plain]
    \centering
    \vfill
    {\usebeamerfont{title}\color{white}\Large #1\par}
    \vfill
  \end{frame}
  \endgroup
}

% and compiling with TEXINPUTS=../Preambles:$TEXINPUTS (the /compile-latex
% skill does this automatically).
%
% PALETTE CONTRACT. Color names defined here MUST match the names in
% Quarto/theme-template.scss so Beamer and Quarto renderings use the same
% palette. The scripts/check-palette-sync.sh script enforces this.
%
% To customize: edit the HEX values below AND the matching $variable in
% Quarto/theme-template.scss, then run scripts/check-palette-sync.sh.
%
% TITLE PAGE. A formal, serif (Libertine) title-page template is defined in
% the Beamer-only block below (\setbeamertemplate{title page}) -- title,
% subtitle, a thin gold rule, then \author/\institute/\date. Frametitle and
% body text remain on Beamer's sans-serif default. No SCSS counterpart is
% needed for this (font/layout only, no new colors).
% =============================================================================

% -----------------------------------------------------------------------------
% Engine detection + encoding
%
% Our /compile-latex skill uses XeLaTeX, where inputenc is unnecessary and
% can produce encoding warnings. Only load inputenc under pdfTeX.
% -----------------------------------------------------------------------------
\usepackage{iftex}
\ifPDFTeX
  \usepackage[utf8]{inputenc}
\fi

\usepackage{amsmath, amssymb, amsthm}
\usepackage{booktabs}
\usepackage{graphicx}

% -----------------------------------------------------------------------------
% Serif face for the title page (formal/academic look). Linux Libertine is
% XeLaTeX- and pdfTeX-safe and redefines \rmdefault, so \rmfamily picks it up
% everywhere \rmfamily is used explicitly. Body text and frametitles stay on
% Beamer's own sans-serif default (unaffected) -- only the title-page
% \setbeamerfont{...} overrides below opt in to \rmfamily.
% -----------------------------------------------------------------------------
\usepackage{libertine}

% -----------------------------------------------------------------------------
% PALETTE — mirrors Quarto/theme-template.scss variable names.
% Load xcolor and define colors BEFORE anything that references them
% (notably hyperref, hypersetup, and the Beamer color assignments).
% -----------------------------------------------------------------------------
\usepackage{xcolor}

\definecolor{primary-blue}{HTML}{012169}      % headings, accents
\definecolor{primary-gold}{HTML}{B9975B}      % emphasis, borders
\definecolor{highlight-yellow}{HTML}{F2A900}  % markers, alerts
\definecolor{light-bg}{HTML}{E8EDF5}          % title / section backgrounds
\definecolor{jet}{HTML}{1A1A1A}               % body text

% Semantic colors (used in diagrams and annotations)
\definecolor{positive}{HTML}{15803D}          % good / observed / identified
\definecolor{negative}{HTML}{B91C1C}          % bad / problematic / bias
\definecolor{neutral}{HTML}{525252}           % reference / context

% Highlight accents (match .hi-* classes in the SCSS)
\definecolor{hi-slate}{HTML}{314F4F}
\definecolor{hi-green}{HTML}{2E8B57}
\definecolor{hi-red}{HTML}{C41E3A}

% Semantic aliases so skill templates and snippets can write
% \textcolor{accent}{...} without hard-coding "primary-blue".
\colorlet{accent}{primary-blue}
\colorlet{accent2}{primary-gold}

% -----------------------------------------------------------------------------
% Hyperref — loaded AFTER xcolor + palette so link colors resolve.
% -----------------------------------------------------------------------------
\usepackage{hyperref}
\hypersetup{colorlinks=true, linkcolor=primary-blue, urlcolor=primary-blue,
            citecolor=primary-blue, pdfborder={0 0 0}}

% -----------------------------------------------------------------------------
% Course code — set once per deck (after % =============================================================================
% Preambles/header.tex — shared LaTeX/Beamer preamble for this project
%
% Use from a Beamer lecture by prepending:
%     % =============================================================================
% Preambles/header.tex — shared LaTeX/Beamer preamble for this project
%
% Use from a Beamer lecture by prepending:
%     \input{header}
% and compiling with TEXINPUTS=../Preambles:$TEXINPUTS (the /compile-latex
% skill does this automatically).
%
% PALETTE CONTRACT. Color names defined here MUST match the names in
% Quarto/theme-template.scss so Beamer and Quarto renderings use the same
% palette. The scripts/check-palette-sync.sh script enforces this.
%
% To customize: edit the HEX values below AND the matching $variable in
% Quarto/theme-template.scss, then run scripts/check-palette-sync.sh.
%
% TITLE PAGE. A formal, serif (Libertine) title-page template is defined in
% the Beamer-only block below (\setbeamertemplate{title page}) -- title,
% subtitle, a thin gold rule, then \author/\institute/\date. Frametitle and
% body text remain on Beamer's sans-serif default. No SCSS counterpart is
% needed for this (font/layout only, no new colors).
% =============================================================================

% -----------------------------------------------------------------------------
% Engine detection + encoding
%
% Our /compile-latex skill uses XeLaTeX, where inputenc is unnecessary and
% can produce encoding warnings. Only load inputenc under pdfTeX.
% -----------------------------------------------------------------------------
\usepackage{iftex}
\ifPDFTeX
  \usepackage[utf8]{inputenc}
\fi

\usepackage{amsmath, amssymb, amsthm}
\usepackage{booktabs}
\usepackage{graphicx}

% -----------------------------------------------------------------------------
% Serif face for the title page (formal/academic look). Linux Libertine is
% XeLaTeX- and pdfTeX-safe and redefines \rmdefault, so \rmfamily picks it up
% everywhere \rmfamily is used explicitly. Body text and frametitles stay on
% Beamer's own sans-serif default (unaffected) -- only the title-page
% \setbeamerfont{...} overrides below opt in to \rmfamily.
% -----------------------------------------------------------------------------
\usepackage{libertine}

% -----------------------------------------------------------------------------
% PALETTE — mirrors Quarto/theme-template.scss variable names.
% Load xcolor and define colors BEFORE anything that references them
% (notably hyperref, hypersetup, and the Beamer color assignments).
% -----------------------------------------------------------------------------
\usepackage{xcolor}

\definecolor{primary-blue}{HTML}{012169}      % headings, accents
\definecolor{primary-gold}{HTML}{B9975B}      % emphasis, borders
\definecolor{highlight-yellow}{HTML}{F2A900}  % markers, alerts
\definecolor{light-bg}{HTML}{E8EDF5}          % title / section backgrounds
\definecolor{jet}{HTML}{1A1A1A}               % body text

% Semantic colors (used in diagrams and annotations)
\definecolor{positive}{HTML}{15803D}          % good / observed / identified
\definecolor{negative}{HTML}{B91C1C}          % bad / problematic / bias
\definecolor{neutral}{HTML}{525252}           % reference / context

% Highlight accents (match .hi-* classes in the SCSS)
\definecolor{hi-slate}{HTML}{314F4F}
\definecolor{hi-green}{HTML}{2E8B57}
\definecolor{hi-red}{HTML}{C41E3A}

% Semantic aliases so skill templates and snippets can write
% \textcolor{accent}{...} without hard-coding "primary-blue".
\colorlet{accent}{primary-blue}
\colorlet{accent2}{primary-gold}

% -----------------------------------------------------------------------------
% Hyperref — loaded AFTER xcolor + palette so link colors resolve.
% -----------------------------------------------------------------------------
\usepackage{hyperref}
\hypersetup{colorlinks=true, linkcolor=primary-blue, urlcolor=primary-blue,
            citecolor=primary-blue, pdfborder={0 0 0}}

% -----------------------------------------------------------------------------
% Course code — set once per deck (after \input{header}, before \begin{document})
% via \coursecode{CS401}. Shown in the footer on every slide so a deck's course
% is identifiable without opening the title page. Empty by default.
% -----------------------------------------------------------------------------
\makeatletter
\newcommand{\coursecode}[1]{\gdef\@coursecodetext{#1}}
\gdef\@coursecodetext{}
\makeatother

% -----------------------------------------------------------------------------
% Beamer theme assignments (only apply under Beamer).
%
% We detect Beamer via \@ifclassloaded{beamer}, which is the documented,
% non-internal way. This must be wrapped in \makeatletter/\makeatother
% because \@ifclassloaded contains an @-catcode character.
% -----------------------------------------------------------------------------
\makeatletter
\@ifclassloaded{beamer}{%
  \usefonttheme{professionalfonts}
  \setbeamercolor{structure}{fg=primary-blue}
  \setbeamercolor{frametitle}{fg=primary-blue}
  \setbeamercolor{title}{fg=primary-blue}
  \setbeamercolor{subtitle}{fg=primary-gold}
  \setbeamercolor{author}{fg=primary-blue}
  \setbeamercolor{institute}{fg=neutral}
  \setbeamercolor{date}{fg=primary-gold}
  \setbeamercolor{itemize item}{fg=primary-blue}
  \setbeamercolor{itemize subitem}{fg=primary-gold}
  \setbeamercolor{alerted text}{fg=primary-gold}
  \setbeamercolor{block title}{fg=white, bg=primary-blue}
  \setbeamercolor{block body}{bg=light-bg}
  % Semantic block variants: block=definition/concept (blue), exampleblock=
  % worked example (green/positive), alertblock=key takeaway or caution (gold).
  % Keep to <=2 colored boxes per slide across ALL three variants (INV-7).
  \setbeamercolor{block title example}{fg=white, bg=positive}
  \setbeamercolor{block body example}{bg=light-bg}
  \setbeamercolor{block title alerted}{fg=white, bg=primary-gold}
  \setbeamercolor{block body alerted}{bg=light-bg}

  % ---------------------------------------------------------------------------
  % Formal title page: serif (Libertine, via \rmfamily) for title-page
  % elements only. Body text, itemize, and frametitle are intentionally left
  % on Beamer's sans-serif default -- \usebeamerfont{title} is also what
  % \transitionslide (below) uses for its big centered text, so section
  % transitions pick up the same serif treatment for free.
  % ---------------------------------------------------------------------------
  \setbeamerfont{title}{family=\rmfamily, series=\bfseries, size=\Large}
  \setbeamerfont{subtitle}{family=\rmfamily, shape=\itshape, size=\normalsize}
  \setbeamerfont{author}{family=\rmfamily, size=\normalsize}
  \setbeamerfont{institute}{family=\rmfamily, size=\scriptsize}
  \setbeamerfont{date}{family=\rmfamily, size=\scriptsize}

  \setbeamertemplate{title page}{
    \vfill
    \begin{center}
      {\usebeamerfont{title}\usebeamercolor[fg]{title}\inserttitle\par}
      \vspace{0.15cm}
      {\usebeamerfont{subtitle}\usebeamercolor[fg]{subtitle}\insertsubtitle\par}
      \vspace{0.45cm}
      {\color{primary-gold}\rule{0.42\textwidth}{0.6pt}\par}
      \vspace{0.45cm}
      {\usebeamerfont{author}\usebeamercolor[fg]{author}\insertauthor\par}
      \vspace{0.35cm}
      {\usebeamerfont{institute}\usebeamercolor[fg]{institute}\insertinstitute\par}
      \vspace{0.35cm}
      {\usebeamerfont{date}\usebeamercolor[fg]{date}\insertdate\par}
    \end{center}
    \vfill
  }

  % Minimal footer: course code (left) + page X / N (right), no navigation clutter
  \setbeamertemplate{navigation symbols}{}
  \setbeamertemplate{footline}{%
    \color{neutral}\footnotesize\hspace{0.7em}\@coursecodetext\hfill
    \insertframenumber/\inserttotalframenumber\hspace{0.7em}\vspace{0.3em}}
}{}
\makeatother

% -----------------------------------------------------------------------------
% TikZ setup — libraries the snippet gallery uses
% -----------------------------------------------------------------------------
\usepackage{tikz}
\usetikzlibrary{arrows.meta, positioning, calc, decorations.pathreplacing,
                fit, shapes.geometric, backgrounds}

% Shared TikZ styles — snippets and hand-written diagrams can pick these up
% via `\begin{tikzpicture}[every node/.style={...}]` or by naming specific
% styles (e.g., `\node[dag-node] (X) at (0,0) {X};`).
\tikzset{
  % Node styles for causal DAGs and similar
  dag-node/.style = {
    draw=primary-blue, fill=light-bg, rounded corners=2pt,
    minimum width=1.2cm, minimum height=0.9cm, align=center, font=\small
  },
  decision-node/.style = {
    draw=primary-gold, fill=white, diamond, aspect=1.8,
    minimum width=1.8cm, minimum height=1.2cm, align=center, font=\small,
    inner sep=1pt
  },
  % Edge styles — observed vs counterfactual
  observed-edge/.style = {-{Latex[length=2.5mm]}, thick, draw=jet},
  counterfactual-edge/.style = {-{Latex[length=2.5mm]}, thick, draw=neutral, dashed},
  confound-edge/.style = {-{Latex[length=2.5mm]}, thick, draw=negative, dashed},
  % Data-point styles for plots
  observed-dot/.style = {fill=primary-blue, draw=primary-blue, circle, inner sep=1.2pt},
  counterfactual-dot/.style = {fill=white, draw=neutral, circle, inner sep=1.2pt}
}

% -----------------------------------------------------------------------------
% Convenience macros
% -----------------------------------------------------------------------------
\newcommand{\muted}[1]{\textcolor{neutral}{#1}}
\newcommand{\key}[1]{\textcolor{primary-gold}{\textbf{#1}}}
\newcommand{\good}[1]{\textcolor{positive}{#1}}
\newcommand{\bad}[1]{\textcolor{negative}{#1}}

% Standout frame macro: full-bleed dark background for section transitions.
% Beamer-only (uses \begin{frame}/\setbeamercolor/\usebeamerfont) -- do not
% call from an article-class file (e.g. Notes/<CODE>/*-notes.tex). \muted,
% \key, \good, \bad, and \coursecode above are plain xcolor/gdef and are
% safe to use from either class; verified when Notes/article-class support
% was added.
% Expand: \transitionslide{Your transition title}
\newcommand{\transitionslide}[1]{%
  \begingroup
  \setbeamercolor{background canvas}{bg=primary-blue}
  \begin{frame}[plain]
    \centering
    \vfill
    {\usebeamerfont{title}\color{white}\Large #1\par}
    \vfill
  \end{frame}
  \endgroup
}

% and compiling with TEXINPUTS=../Preambles:$TEXINPUTS (the /compile-latex
% skill does this automatically).
%
% PALETTE CONTRACT. Color names defined here MUST match the names in
% Quarto/theme-template.scss so Beamer and Quarto renderings use the same
% palette. The scripts/check-palette-sync.sh script enforces this.
%
% To customize: edit the HEX values below AND the matching $variable in
% Quarto/theme-template.scss, then run scripts/check-palette-sync.sh.
%
% TITLE PAGE. A formal, serif (Libertine) title-page template is defined in
% the Beamer-only block below (\setbeamertemplate{title page}) -- title,
% subtitle, a thin gold rule, then \author/\institute/\date. Frametitle and
% body text remain on Beamer's sans-serif default. No SCSS counterpart is
% needed for this (font/layout only, no new colors).
% =============================================================================

% -----------------------------------------------------------------------------
% Engine detection + encoding
%
% Our /compile-latex skill uses XeLaTeX, where inputenc is unnecessary and
% can produce encoding warnings. Only load inputenc under pdfTeX.
% -----------------------------------------------------------------------------
\usepackage{iftex}
\ifPDFTeX
  \usepackage[utf8]{inputenc}
\fi

\usepackage{amsmath, amssymb, amsthm}
\usepackage{booktabs}
\usepackage{graphicx}

% -----------------------------------------------------------------------------
% Serif face for the title page (formal/academic look). Linux Libertine is
% XeLaTeX- and pdfTeX-safe and redefines \rmdefault, so \rmfamily picks it up
% everywhere \rmfamily is used explicitly. Body text and frametitles stay on
% Beamer's own sans-serif default (unaffected) -- only the title-page
% \setbeamerfont{...} overrides below opt in to \rmfamily.
% -----------------------------------------------------------------------------
\usepackage{libertine}

% -----------------------------------------------------------------------------
% PALETTE — mirrors Quarto/theme-template.scss variable names.
% Load xcolor and define colors BEFORE anything that references them
% (notably hyperref, hypersetup, and the Beamer color assignments).
% -----------------------------------------------------------------------------
\usepackage{xcolor}

\definecolor{primary-blue}{HTML}{012169}      % headings, accents
\definecolor{primary-gold}{HTML}{B9975B}      % emphasis, borders
\definecolor{highlight-yellow}{HTML}{F2A900}  % markers, alerts
\definecolor{light-bg}{HTML}{E8EDF5}          % title / section backgrounds
\definecolor{jet}{HTML}{1A1A1A}               % body text

% Semantic colors (used in diagrams and annotations)
\definecolor{positive}{HTML}{15803D}          % good / observed / identified
\definecolor{negative}{HTML}{B91C1C}          % bad / problematic / bias
\definecolor{neutral}{HTML}{525252}           % reference / context

% Highlight accents (match .hi-* classes in the SCSS)
\definecolor{hi-slate}{HTML}{314F4F}
\definecolor{hi-green}{HTML}{2E8B57}
\definecolor{hi-red}{HTML}{C41E3A}

% Semantic aliases so skill templates and snippets can write
% \textcolor{accent}{...} without hard-coding "primary-blue".
\colorlet{accent}{primary-blue}
\colorlet{accent2}{primary-gold}

% -----------------------------------------------------------------------------
% Hyperref — loaded AFTER xcolor + palette so link colors resolve.
% -----------------------------------------------------------------------------
\usepackage{hyperref}
\hypersetup{colorlinks=true, linkcolor=primary-blue, urlcolor=primary-blue,
            citecolor=primary-blue, pdfborder={0 0 0}}

% -----------------------------------------------------------------------------
% Course code — set once per deck (after % =============================================================================
% Preambles/header.tex — shared LaTeX/Beamer preamble for this project
%
% Use from a Beamer lecture by prepending:
%     \input{header}
% and compiling with TEXINPUTS=../Preambles:$TEXINPUTS (the /compile-latex
% skill does this automatically).
%
% PALETTE CONTRACT. Color names defined here MUST match the names in
% Quarto/theme-template.scss so Beamer and Quarto renderings use the same
% palette. The scripts/check-palette-sync.sh script enforces this.
%
% To customize: edit the HEX values below AND the matching $variable in
% Quarto/theme-template.scss, then run scripts/check-palette-sync.sh.
%
% TITLE PAGE. A formal, serif (Libertine) title-page template is defined in
% the Beamer-only block below (\setbeamertemplate{title page}) -- title,
% subtitle, a thin gold rule, then \author/\institute/\date. Frametitle and
% body text remain on Beamer's sans-serif default. No SCSS counterpart is
% needed for this (font/layout only, no new colors).
% =============================================================================

% -----------------------------------------------------------------------------
% Engine detection + encoding
%
% Our /compile-latex skill uses XeLaTeX, where inputenc is unnecessary and
% can produce encoding warnings. Only load inputenc under pdfTeX.
% -----------------------------------------------------------------------------
\usepackage{iftex}
\ifPDFTeX
  \usepackage[utf8]{inputenc}
\fi

\usepackage{amsmath, amssymb, amsthm}
\usepackage{booktabs}
\usepackage{graphicx}

% -----------------------------------------------------------------------------
% Serif face for the title page (formal/academic look). Linux Libertine is
% XeLaTeX- and pdfTeX-safe and redefines \rmdefault, so \rmfamily picks it up
% everywhere \rmfamily is used explicitly. Body text and frametitles stay on
% Beamer's own sans-serif default (unaffected) -- only the title-page
% \setbeamerfont{...} overrides below opt in to \rmfamily.
% -----------------------------------------------------------------------------
\usepackage{libertine}

% -----------------------------------------------------------------------------
% PALETTE — mirrors Quarto/theme-template.scss variable names.
% Load xcolor and define colors BEFORE anything that references them
% (notably hyperref, hypersetup, and the Beamer color assignments).
% -----------------------------------------------------------------------------
\usepackage{xcolor}

\definecolor{primary-blue}{HTML}{012169}      % headings, accents
\definecolor{primary-gold}{HTML}{B9975B}      % emphasis, borders
\definecolor{highlight-yellow}{HTML}{F2A900}  % markers, alerts
\definecolor{light-bg}{HTML}{E8EDF5}          % title / section backgrounds
\definecolor{jet}{HTML}{1A1A1A}               % body text

% Semantic colors (used in diagrams and annotations)
\definecolor{positive}{HTML}{15803D}          % good / observed / identified
\definecolor{negative}{HTML}{B91C1C}          % bad / problematic / bias
\definecolor{neutral}{HTML}{525252}           % reference / context

% Highlight accents (match .hi-* classes in the SCSS)
\definecolor{hi-slate}{HTML}{314F4F}
\definecolor{hi-green}{HTML}{2E8B57}
\definecolor{hi-red}{HTML}{C41E3A}

% Semantic aliases so skill templates and snippets can write
% \textcolor{accent}{...} without hard-coding "primary-blue".
\colorlet{accent}{primary-blue}
\colorlet{accent2}{primary-gold}

% -----------------------------------------------------------------------------
% Hyperref — loaded AFTER xcolor + palette so link colors resolve.
% -----------------------------------------------------------------------------
\usepackage{hyperref}
\hypersetup{colorlinks=true, linkcolor=primary-blue, urlcolor=primary-blue,
            citecolor=primary-blue, pdfborder={0 0 0}}

% -----------------------------------------------------------------------------
% Course code — set once per deck (after \input{header}, before \begin{document})
% via \coursecode{CS401}. Shown in the footer on every slide so a deck's course
% is identifiable without opening the title page. Empty by default.
% -----------------------------------------------------------------------------
\makeatletter
\newcommand{\coursecode}[1]{\gdef\@coursecodetext{#1}}
\gdef\@coursecodetext{}
\makeatother

% -----------------------------------------------------------------------------
% Beamer theme assignments (only apply under Beamer).
%
% We detect Beamer via \@ifclassloaded{beamer}, which is the documented,
% non-internal way. This must be wrapped in \makeatletter/\makeatother
% because \@ifclassloaded contains an @-catcode character.
% -----------------------------------------------------------------------------
\makeatletter
\@ifclassloaded{beamer}{%
  \usefonttheme{professionalfonts}
  \setbeamercolor{structure}{fg=primary-blue}
  \setbeamercolor{frametitle}{fg=primary-blue}
  \setbeamercolor{title}{fg=primary-blue}
  \setbeamercolor{subtitle}{fg=primary-gold}
  \setbeamercolor{author}{fg=primary-blue}
  \setbeamercolor{institute}{fg=neutral}
  \setbeamercolor{date}{fg=primary-gold}
  \setbeamercolor{itemize item}{fg=primary-blue}
  \setbeamercolor{itemize subitem}{fg=primary-gold}
  \setbeamercolor{alerted text}{fg=primary-gold}
  \setbeamercolor{block title}{fg=white, bg=primary-blue}
  \setbeamercolor{block body}{bg=light-bg}
  % Semantic block variants: block=definition/concept (blue), exampleblock=
  % worked example (green/positive), alertblock=key takeaway or caution (gold).
  % Keep to <=2 colored boxes per slide across ALL three variants (INV-7).
  \setbeamercolor{block title example}{fg=white, bg=positive}
  \setbeamercolor{block body example}{bg=light-bg}
  \setbeamercolor{block title alerted}{fg=white, bg=primary-gold}
  \setbeamercolor{block body alerted}{bg=light-bg}

  % ---------------------------------------------------------------------------
  % Formal title page: serif (Libertine, via \rmfamily) for title-page
  % elements only. Body text, itemize, and frametitle are intentionally left
  % on Beamer's sans-serif default -- \usebeamerfont{title} is also what
  % \transitionslide (below) uses for its big centered text, so section
  % transitions pick up the same serif treatment for free.
  % ---------------------------------------------------------------------------
  \setbeamerfont{title}{family=\rmfamily, series=\bfseries, size=\Large}
  \setbeamerfont{subtitle}{family=\rmfamily, shape=\itshape, size=\normalsize}
  \setbeamerfont{author}{family=\rmfamily, size=\normalsize}
  \setbeamerfont{institute}{family=\rmfamily, size=\scriptsize}
  \setbeamerfont{date}{family=\rmfamily, size=\scriptsize}

  \setbeamertemplate{title page}{
    \vfill
    \begin{center}
      {\usebeamerfont{title}\usebeamercolor[fg]{title}\inserttitle\par}
      \vspace{0.15cm}
      {\usebeamerfont{subtitle}\usebeamercolor[fg]{subtitle}\insertsubtitle\par}
      \vspace{0.45cm}
      {\color{primary-gold}\rule{0.42\textwidth}{0.6pt}\par}
      \vspace{0.45cm}
      {\usebeamerfont{author}\usebeamercolor[fg]{author}\insertauthor\par}
      \vspace{0.35cm}
      {\usebeamerfont{institute}\usebeamercolor[fg]{institute}\insertinstitute\par}
      \vspace{0.35cm}
      {\usebeamerfont{date}\usebeamercolor[fg]{date}\insertdate\par}
    \end{center}
    \vfill
  }

  % Minimal footer: course code (left) + page X / N (right), no navigation clutter
  \setbeamertemplate{navigation symbols}{}
  \setbeamertemplate{footline}{%
    \color{neutral}\footnotesize\hspace{0.7em}\@coursecodetext\hfill
    \insertframenumber/\inserttotalframenumber\hspace{0.7em}\vspace{0.3em}}
}{}
\makeatother

% -----------------------------------------------------------------------------
% TikZ setup — libraries the snippet gallery uses
% -----------------------------------------------------------------------------
\usepackage{tikz}
\usetikzlibrary{arrows.meta, positioning, calc, decorations.pathreplacing,
                fit, shapes.geometric, backgrounds}

% Shared TikZ styles — snippets and hand-written diagrams can pick these up
% via `\begin{tikzpicture}[every node/.style={...}]` or by naming specific
% styles (e.g., `\node[dag-node] (X) at (0,0) {X};`).
\tikzset{
  % Node styles for causal DAGs and similar
  dag-node/.style = {
    draw=primary-blue, fill=light-bg, rounded corners=2pt,
    minimum width=1.2cm, minimum height=0.9cm, align=center, font=\small
  },
  decision-node/.style = {
    draw=primary-gold, fill=white, diamond, aspect=1.8,
    minimum width=1.8cm, minimum height=1.2cm, align=center, font=\small,
    inner sep=1pt
  },
  % Edge styles — observed vs counterfactual
  observed-edge/.style = {-{Latex[length=2.5mm]}, thick, draw=jet},
  counterfactual-edge/.style = {-{Latex[length=2.5mm]}, thick, draw=neutral, dashed},
  confound-edge/.style = {-{Latex[length=2.5mm]}, thick, draw=negative, dashed},
  % Data-point styles for plots
  observed-dot/.style = {fill=primary-blue, draw=primary-blue, circle, inner sep=1.2pt},
  counterfactual-dot/.style = {fill=white, draw=neutral, circle, inner sep=1.2pt}
}

% -----------------------------------------------------------------------------
% Convenience macros
% -----------------------------------------------------------------------------
\newcommand{\muted}[1]{\textcolor{neutral}{#1}}
\newcommand{\key}[1]{\textcolor{primary-gold}{\textbf{#1}}}
\newcommand{\good}[1]{\textcolor{positive}{#1}}
\newcommand{\bad}[1]{\textcolor{negative}{#1}}

% Standout frame macro: full-bleed dark background for section transitions.
% Beamer-only (uses \begin{frame}/\setbeamercolor/\usebeamerfont) -- do not
% call from an article-class file (e.g. Notes/<CODE>/*-notes.tex). \muted,
% \key, \good, \bad, and \coursecode above are plain xcolor/gdef and are
% safe to use from either class; verified when Notes/article-class support
% was added.
% Expand: \transitionslide{Your transition title}
\newcommand{\transitionslide}[1]{%
  \begingroup
  \setbeamercolor{background canvas}{bg=primary-blue}
  \begin{frame}[plain]
    \centering
    \vfill
    {\usebeamerfont{title}\color{white}\Large #1\par}
    \vfill
  \end{frame}
  \endgroup
}
, before \begin{document})
% via \coursecode{CS401}. Shown in the footer on every slide so a deck's course
% is identifiable without opening the title page. Empty by default.
% -----------------------------------------------------------------------------
\makeatletter
\newcommand{\coursecode}[1]{\gdef\@coursecodetext{#1}}
\gdef\@coursecodetext{}
\makeatother

% -----------------------------------------------------------------------------
% Beamer theme assignments (only apply under Beamer).
%
% We detect Beamer via \@ifclassloaded{beamer}, which is the documented,
% non-internal way. This must be wrapped in \makeatletter/\makeatother
% because \@ifclassloaded contains an @-catcode character.
% -----------------------------------------------------------------------------
\makeatletter
\@ifclassloaded{beamer}{%
  \usefonttheme{professionalfonts}
  \setbeamercolor{structure}{fg=primary-blue}
  \setbeamercolor{frametitle}{fg=primary-blue}
  \setbeamercolor{title}{fg=primary-blue}
  \setbeamercolor{subtitle}{fg=primary-gold}
  \setbeamercolor{author}{fg=primary-blue}
  \setbeamercolor{institute}{fg=neutral}
  \setbeamercolor{date}{fg=primary-gold}
  \setbeamercolor{itemize item}{fg=primary-blue}
  \setbeamercolor{itemize subitem}{fg=primary-gold}
  \setbeamercolor{alerted text}{fg=primary-gold}
  \setbeamercolor{block title}{fg=white, bg=primary-blue}
  \setbeamercolor{block body}{bg=light-bg}
  % Semantic block variants: block=definition/concept (blue), exampleblock=
  % worked example (green/positive), alertblock=key takeaway or caution (gold).
  % Keep to <=2 colored boxes per slide across ALL three variants (INV-7).
  \setbeamercolor{block title example}{fg=white, bg=positive}
  \setbeamercolor{block body example}{bg=light-bg}
  \setbeamercolor{block title alerted}{fg=white, bg=primary-gold}
  \setbeamercolor{block body alerted}{bg=light-bg}

  % ---------------------------------------------------------------------------
  % Formal title page: serif (Libertine, via \rmfamily) for title-page
  % elements only. Body text, itemize, and frametitle are intentionally left
  % on Beamer's sans-serif default -- \usebeamerfont{title} is also what
  % \transitionslide (below) uses for its big centered text, so section
  % transitions pick up the same serif treatment for free.
  % ---------------------------------------------------------------------------
  \setbeamerfont{title}{family=\rmfamily, series=\bfseries, size=\Large}
  \setbeamerfont{subtitle}{family=\rmfamily, shape=\itshape, size=\normalsize}
  \setbeamerfont{author}{family=\rmfamily, size=\normalsize}
  \setbeamerfont{institute}{family=\rmfamily, size=\scriptsize}
  \setbeamerfont{date}{family=\rmfamily, size=\scriptsize}

  \setbeamertemplate{title page}{
    \vfill
    \begin{center}
      {\usebeamerfont{title}\usebeamercolor[fg]{title}\inserttitle\par}
      \vspace{0.15cm}
      {\usebeamerfont{subtitle}\usebeamercolor[fg]{subtitle}\insertsubtitle\par}
      \vspace{0.45cm}
      {\color{primary-gold}\rule{0.42\textwidth}{0.6pt}\par}
      \vspace{0.45cm}
      {\usebeamerfont{author}\usebeamercolor[fg]{author}\insertauthor\par}
      \vspace{0.35cm}
      {\usebeamerfont{institute}\usebeamercolor[fg]{institute}\insertinstitute\par}
      \vspace{0.35cm}
      {\usebeamerfont{date}\usebeamercolor[fg]{date}\insertdate\par}
    \end{center}
    \vfill
  }

  % Minimal footer: course code (left) + page X / N (right), no navigation clutter
  \setbeamertemplate{navigation symbols}{}
  \setbeamertemplate{footline}{%
    \color{neutral}\footnotesize\hspace{0.7em}\@coursecodetext\hfill
    \insertframenumber/\inserttotalframenumber\hspace{0.7em}\vspace{0.3em}}
}{}
\makeatother

% -----------------------------------------------------------------------------
% TikZ setup — libraries the snippet gallery uses
% -----------------------------------------------------------------------------
\usepackage{tikz}
\usetikzlibrary{arrows.meta, positioning, calc, decorations.pathreplacing,
                fit, shapes.geometric, backgrounds}

% Shared TikZ styles — snippets and hand-written diagrams can pick these up
% via `\begin{tikzpicture}[every node/.style={...}]` or by naming specific
% styles (e.g., `\node[dag-node] (X) at (0,0) {X};`).
\tikzset{
  % Node styles for causal DAGs and similar
  dag-node/.style = {
    draw=primary-blue, fill=light-bg, rounded corners=2pt,
    minimum width=1.2cm, minimum height=0.9cm, align=center, font=\small
  },
  decision-node/.style = {
    draw=primary-gold, fill=white, diamond, aspect=1.8,
    minimum width=1.8cm, minimum height=1.2cm, align=center, font=\small,
    inner sep=1pt
  },
  % Edge styles — observed vs counterfactual
  observed-edge/.style = {-{Latex[length=2.5mm]}, thick, draw=jet},
  counterfactual-edge/.style = {-{Latex[length=2.5mm]}, thick, draw=neutral, dashed},
  confound-edge/.style = {-{Latex[length=2.5mm]}, thick, draw=negative, dashed},
  % Data-point styles for plots
  observed-dot/.style = {fill=primary-blue, draw=primary-blue, circle, inner sep=1.2pt},
  counterfactual-dot/.style = {fill=white, draw=neutral, circle, inner sep=1.2pt}
}

% -----------------------------------------------------------------------------
% Convenience macros
% -----------------------------------------------------------------------------
\newcommand{\muted}[1]{\textcolor{neutral}{#1}}
\newcommand{\key}[1]{\textcolor{primary-gold}{\textbf{#1}}}
\newcommand{\good}[1]{\textcolor{positive}{#1}}
\newcommand{\bad}[1]{\textcolor{negative}{#1}}

% Standout frame macro: full-bleed dark background for section transitions.
% Beamer-only (uses \begin{frame}/\setbeamercolor/\usebeamerfont) -- do not
% call from an article-class file (e.g. Notes/<CODE>/*-notes.tex). \muted,
% \key, \good, \bad, and \coursecode above are plain xcolor/gdef and are
% safe to use from either class; verified when Notes/article-class support
% was added.
% Expand: \transitionslide{Your transition title}
\newcommand{\transitionslide}[1]{%
  \begingroup
  \setbeamercolor{background canvas}{bg=primary-blue}
  \begin{frame}[plain]
    \centering
    \vfill
    {\usebeamerfont{title}\color{white}\Large #1\par}
    \vfill
  \end{frame}
  \endgroup
}
, before \begin{document})
% via \coursecode{CS401}. Shown in the footer on every slide so a deck's course
% is identifiable without opening the title page. Empty by default.
% -----------------------------------------------------------------------------
\makeatletter
\newcommand{\coursecode}[1]{\gdef\@coursecodetext{#1}}
\gdef\@coursecodetext{}
\makeatother

% -----------------------------------------------------------------------------
% Beamer theme assignments (only apply under Beamer).
%
% We detect Beamer via \@ifclassloaded{beamer}, which is the documented,
% non-internal way. This must be wrapped in \makeatletter/\makeatother
% because \@ifclassloaded contains an @-catcode character.
% -----------------------------------------------------------------------------
\makeatletter
\@ifclassloaded{beamer}{%
  \usefonttheme{professionalfonts}
  \setbeamercolor{structure}{fg=primary-blue}
  \setbeamercolor{frametitle}{fg=primary-blue}
  \setbeamercolor{title}{fg=primary-blue}
  \setbeamercolor{subtitle}{fg=primary-gold}
  \setbeamercolor{author}{fg=primary-blue}
  \setbeamercolor{institute}{fg=neutral}
  \setbeamercolor{date}{fg=primary-gold}
  \setbeamercolor{itemize item}{fg=primary-blue}
  \setbeamercolor{itemize subitem}{fg=primary-gold}
  \setbeamercolor{alerted text}{fg=primary-gold}
  \setbeamercolor{block title}{fg=white, bg=primary-blue}
  \setbeamercolor{block body}{bg=light-bg}
  % Semantic block variants: block=definition/concept (blue), exampleblock=
  % worked example (green/positive), alertblock=key takeaway or caution (gold).
  % Keep to <=2 colored boxes per slide across ALL three variants (INV-7).
  \setbeamercolor{block title example}{fg=white, bg=positive}
  \setbeamercolor{block body example}{bg=light-bg}
  \setbeamercolor{block title alerted}{fg=white, bg=primary-gold}
  \setbeamercolor{block body alerted}{bg=light-bg}

  % ---------------------------------------------------------------------------
  % Formal title page: serif (Libertine, via \rmfamily) for title-page
  % elements only. Body text, itemize, and frametitle are intentionally left
  % on Beamer's sans-serif default -- \usebeamerfont{title} is also what
  % \transitionslide (below) uses for its big centered text, so section
  % transitions pick up the same serif treatment for free.
  % ---------------------------------------------------------------------------
  \setbeamerfont{title}{family=\rmfamily, series=\bfseries, size=\Large}
  \setbeamerfont{subtitle}{family=\rmfamily, shape=\itshape, size=\normalsize}
  \setbeamerfont{author}{family=\rmfamily, size=\normalsize}
  \setbeamerfont{institute}{family=\rmfamily, size=\scriptsize}
  \setbeamerfont{date}{family=\rmfamily, size=\scriptsize}

  \setbeamertemplate{title page}{
    \vfill
    \begin{center}
      {\usebeamerfont{title}\usebeamercolor[fg]{title}\inserttitle\par}
      \vspace{0.15cm}
      {\usebeamerfont{subtitle}\usebeamercolor[fg]{subtitle}\insertsubtitle\par}
      \vspace{0.45cm}
      {\color{primary-gold}\rule{0.42\textwidth}{0.6pt}\par}
      \vspace{0.45cm}
      {\usebeamerfont{author}\usebeamercolor[fg]{author}\insertauthor\par}
      \vspace{0.35cm}
      {\usebeamerfont{institute}\usebeamercolor[fg]{institute}\insertinstitute\par}
      \vspace{0.35cm}
      {\usebeamerfont{date}\usebeamercolor[fg]{date}\insertdate\par}
    \end{center}
    \vfill
  }

  % Minimal footer: course code (left) + page X / N (right), no navigation clutter
  \setbeamertemplate{navigation symbols}{}
  \setbeamertemplate{footline}{%
    \color{neutral}\footnotesize\hspace{0.7em}\@coursecodetext\hfill
    \insertframenumber/\inserttotalframenumber\hspace{0.7em}\vspace{0.3em}}
}{}
\makeatother

% -----------------------------------------------------------------------------
% TikZ setup — libraries the snippet gallery uses
% -----------------------------------------------------------------------------
\usepackage{tikz}
\usetikzlibrary{arrows.meta, positioning, calc, decorations.pathreplacing,
                fit, shapes.geometric, backgrounds}

% Shared TikZ styles — snippets and hand-written diagrams can pick these up
% via `\begin{tikzpicture}[every node/.style={...}]` or by naming specific
% styles (e.g., `\node[dag-node] (X) at (0,0) {X};`).
\tikzset{
  % Node styles for causal DAGs and similar
  dag-node/.style = {
    draw=primary-blue, fill=light-bg, rounded corners=2pt,
    minimum width=1.2cm, minimum height=0.9cm, align=center, font=\small
  },
  decision-node/.style = {
    draw=primary-gold, fill=white, diamond, aspect=1.8,
    minimum width=1.8cm, minimum height=1.2cm, align=center, font=\small,
    inner sep=1pt
  },
  % Edge styles — observed vs counterfactual
  observed-edge/.style = {-{Latex[length=2.5mm]}, thick, draw=jet},
  counterfactual-edge/.style = {-{Latex[length=2.5mm]}, thick, draw=neutral, dashed},
  confound-edge/.style = {-{Latex[length=2.5mm]}, thick, draw=negative, dashed},
  % Data-point styles for plots
  observed-dot/.style = {fill=primary-blue, draw=primary-blue, circle, inner sep=1.2pt},
  counterfactual-dot/.style = {fill=white, draw=neutral, circle, inner sep=1.2pt}
}

% -----------------------------------------------------------------------------
% Convenience macros
% -----------------------------------------------------------------------------
\newcommand{\muted}[1]{\textcolor{neutral}{#1}}
\newcommand{\key}[1]{\textcolor{primary-gold}{\textbf{#1}}}
\newcommand{\good}[1]{\textcolor{positive}{#1}}
\newcommand{\bad}[1]{\textcolor{negative}{#1}}

% Standout frame macro: full-bleed dark background for section transitions.
% Beamer-only (uses \begin{frame}/\setbeamercolor/\usebeamerfont) -- do not
% call from an article-class file (e.g. Notes/<CODE>/*-notes.tex). \muted,
% \key, \good, \bad, and \coursecode above are plain xcolor/gdef and are
% safe to use from either class; verified when Notes/article-class support
% was added.
% Expand: \transitionslide{Your transition title}
\newcommand{\transitionslide}[1]{%
  \begingroup
  \setbeamercolor{background canvas}{bg=primary-blue}
  \begin{frame}[plain]
    \centering
    \vfill
    {\usebeamerfont{title}\color{white}\Large #1\par}
    \vfill
  \end{frame}
  \endgroup
}

\coursecode{CS401}
\renewcommand{\thesection}{4.\arabic{section}}

\title{Competitive Exam Practice 4 --- Answer Key\\\large Registers, Instructions, and the Instruction Cycle}
\author{Auqib Hamid Lone\\[0.3em]
  \emph{Department of Computer Science \& Engineering}, \emph{GCET Ganderbal}}
\date{\today}

\begin{document}
\maketitle

\begin{enumerate}
\item[\textbf{Q1}] \textbf{256.} Expanding-opcode decomposition:
\begin{itemize}
  \item \textbf{Two-address:} opcode occupies $11 - 4 - 4 = 3$ bits $\Rightarrow 2^3 = 8$ codes; 5 used, leaving 3.
  \item \textbf{One-address:} each of the 3 free two-address prefixes extends into the former second-address field: $3 \times 2^4 = 48$ codes; 32 used, leaving $48 - 32 = 16$.
  \item \textbf{Zero-address:} each of the 16 free one-address prefixes extends into the last 4-bit field: $16 \times 2^4 = 256$ codes.
\end{itemize}

\item[\textbf{Q2}] \textbf{16383.} $32$-bit word $=$ opcode $+$ two register fields $+$ immediate. Each register needs $\lceil\log_2 64\rceil = 6$ bits, so the two fields cost 12 bits. The opcode needs $\lceil\log_2 45\rceil = 6$ bits (6 bits encode 64 $\geq$ 45; 5 bits encode only 32). Remaining immediate width: $32 - 6 - 12 = 14$ bits. Largest unsigned $= 2^{14} - 1 = 16383$.

\item[\textbf{Q3}] \textbf{32.} $32$-bit word $=$ opcode $+$ mode $+$ two register fields $+$ scalar. Addressing-mode field: $\lceil\log_2 8\rceil = 3$ bits. Each register: $\lceil\log_2 16\rceil = 4$ bits, so two fields cost 8 bits. Scalar field: 16 bits. Opcode width $= 32 - 3 - 8 - 16 = 5$ bits, giving $2^5 = 32$ unique opcodes per addressing mode.

\item[\textbf{Q4}] \textbf{(B) 3.} On a single-bus data path the execution of \texttt{add R0, R1} needs one value on the bus at a time:
\begin{enumerate}
  \item Drive \texttt{[R1]} onto the bus into one ALU input.
  \item Drive \texttt{[R0]} onto the bus into the other ALU input.
  \item ALU computes $R0 \gets R0 + R1$; result written back to \texttt{R0}.
\end{enumerate}
No memory access is involved (both operands are in registers), so the two-cycle memory-read timing does not affect this instruction --- the minimum is 3 clock cycles.

\item[\textbf{Q5}] \textbf{(C) 4.} \texttt{call Rn, sub} is a two-word instruction, so its execution must fetch the second word and then transfer:
\begin{enumerate}
  \item \texttt{MAR} $\gets$ \texttt{[PC]} --- load the address of the second word (memory read cycle 1).
  \item \texttt{MDR} $\gets$ \emph{memory data} --- fetch the second word (memory read cycle 2).
  \item $Rn \gets PC + 1$ --- ALU computes the return address (saved during fetch of the first word, so \texttt{PC} now points at the second word).
  \item $PC \gets \texttt{[MDR]}$ --- jump target loaded into the PC.
\end{enumerate}
Total: 4 clock cycles.

\item[\textbf{Q6}] \textbf{3.} The XOR-swap uses the identity $x \oplus y$ three times:
\begin{itemize}
  \item $XOR\ R_1, R_2$ \quad ($R_1 \gets R_1 \oplus R_2$)
  \item $XOR\ R_2, R_1$ \quad ($R_2 \gets (R_1 \oplus R_2) \oplus R_2 = R_1$)
  \item $XOR\ R_1, R_2$ \quad ($R_1 \gets (R_1 \oplus R_2) \oplus R_1 = R_2$)
\end{itemize}
After three XORs the contents are exchanged. No extra register or temporary storage is needed, so 3 is minimal.

\item[\textbf{Q7}] \textbf{-8.} The branch is at instruction $i{+}3$; the offset is relative to the \emph{next} instruction, i.e.\ $i{+}4$. The target is $i{+}2$, which is $2$ instructions before the reference point, i.e.\ $-2$ instructions $\times$ 4 bytes per instruction $= -8$ bytes.

\item[\textbf{Q8}] \textbf{(A).} PUSH decrements first so the operand is stored at the \emph{new} (lower) top: $SP \gets [SP]-1$ (stack grows downward toward lower addresses), then $M[SP] \gets$ operand. Writing the operand first would overwrite the current top before \texttt{SP} moves, and (C)/(D) use $+1$, which would grow the stack upward.
\end{enumerate}

\end{document}
